{
\footnotesize
\begin{longtable}{|m{2cm}|m{2.5cm}|m{2.5cm}|m{3cm}|m{2cm}|}
\caption{\raggedright Perbandingan Detail Arsitektur Saat Ini dengan yang Diusulkan}
\label{tab:architecture_comparison} \\
\hline
\multicolumn{1}{|>{\centering\arraybackslash}m{2cm}|}{\textbf{Aspek}} &
\multicolumn{1}{>{\centering\arraybackslash}m{2.5cm}|}{\textbf{Arsitektur Saat Ini}} &
\multicolumn{1}{>{\centering\arraybackslash}m{2.5cm}|}{\textbf{Arsitektur Diusulkan}} &
\multicolumn{1}{>{\centering\arraybackslash}m{3cm}|}{\textbf{Keuntungan Terukur}} &
\multicolumn{1}{>{\centering\arraybackslash}m{2cm}|}{\textbf{Referensi}} \\
\hline
\endfirsthead
\hline
\multicolumn{1}{|>{\centering\arraybackslash}m{2cm}|}{\textbf{Aspek}} &
\multicolumn{1}{>{\centering\arraybackslash}m{2.5cm}|}{\textbf{Arsitektur Saat Ini}} &
\multicolumn{1}{>{\centering\arraybackslash}m{2.5cm}|}{\textbf{Arsitektur Diusulkan}} &
\multicolumn{1}{>{\centering\arraybackslash}m{3cm}|}{\textbf{Keuntungan Terukur}} &
\multicolumn{1}{>{\centering\arraybackslash}m{2cm}|}{\textbf{Referensi}} \\
\hline
\endhead
\hline
\endfoot
\hline
\endlastfoot

Orkestrasi &
\textit{Compute} heterogen dengan \textit{provisioning} manual, waktu tunggu \textit{scaling} panjang, dan tingkat standardisasi rendah. &
Platform Kubernetes terpadu dengan manajemen deklaratif dan alokasi sumber daya otomatis, sejalan dengan komponen arsitektur MLOps yang diidentifikasi \textcite{kreuzberger2023mlops}. GitOps dengan Argo CD memastikan perubahan infrastruktur dan aplikasi mengikuti \textit{single source of truth} di Git. &
\textcite{lakka2025kubernetes} menunjukkan bahwa adopsi Kubernetes pada perusahaan disertai praktik DevOps yang matang meningkatkan \textit{deployment frequency} dan menekan \textit{Mean Time To Recovery} (MTTR) secara signifikan; arsitektur yang diusulkan menargetkan karakteristik performa pada arah yang sama. &
\autocite{kreuzberger2023mlops,lakka2025kubernetes} \\
\hline
Manajemen Fitur &
Logika fitur tersebar di \textit{notebook} dan skrip dengan duplikasi tinggi dan banyak fungsi yang saling tumpang tindih. &
Feast dengan \textit{feature registry} terpusat, \textit{dual serving} \textit{offline/online}, serta dukungan \textit{point-in-time correctness}. &
Duplikasi fitur dapat ditekan, \textit{training--serving skew} berkurang melalui pemisahan \textit{offline/online} dan \textit{point-in-time retrieval}, dan \textit{feature reuse} berpotensi meningkat signifikan. &
\autocite{munappy2022data,polyzotis2018data} \\
\hline
\textit{Vector Embeddings} &
Fitur \textit{embedding} tidak didukung secara \textit{native} atau diimplementasikan secara kustom dengan latensi tinggi dan \textit{throughput} terbatas. &
Redis Stack dengan dukungan HNSW \textit{vector search} berlatensi rendah. &
\textcite{johnson2021billion} menunjukkan bahwa indeks vektor berbasis HNSW maupun FAISS dapat melayani \textit{billion-scale similarity search} dengan latensi rendah; desain ini memanfaatkan prinsip serupa pada Redis Stack agar fitur \textit{embedding} dilayani secara \textit{online} tanpa membangun layanan kustom. &
\citereset\autocite{johnson2021billion} \\
\hline
Validitas Temporal &
\textit{Point-in-time} \textit{join} dilakukan manual dan rentan kesalahan, sehingga banyak tim mengalami \textit{data leakage}. &
\textit{Point-in-time join} otomatis dari Feast dengan pola verifikasi yang lebih sistematis. &
Pendekatan ini mencegah kebocoran data temporal yang dapat menghasilkan estimasi performa \textit{backtest} yang terlalu optimistis dan tidak realistis di produksi. &
\autocite{polyzotis2018data} \\
\hline
Pemrosesan Data &
\textit{Batch} dan \textit{stream} diproses oleh \textit{pipeline} terpisah tanpa arsitektur yang koheren, sehingga hasil pemrosesan tidak konsisten. &
Spark sebagai mesin \textit{batch processing} dan Flink sebagai mesin \textit{stream processing} yang hasilnya dimaterialisasikan langsung ke \textit{Feature Store}. Tekton dan Argo Workflows mengorkestrasi \textit{pipeline} berbasis manifes deklaratif sehingga pengembangan dan penjadwalan transformasi menjadi lebih sistematis. &
Waktu pengembangan \textit{pipeline} dapat menurun dan \textit{freshness} data meningkat berkat pemrosesan \textit{streaming} yang lebih efisien dan terintegrasi. &
\autocite{rodrigues2025architecture} \\
\hline
Pelatihan Model &
\textit{Workflow} pelatihan berbasis \textit{notebook} dengan \textit{tracking} eksperimen informal dan \textit{artifact} tidak terkelola. &
Kubeflow \textit{Pipelines} untuk orkestrasi pelatihan dan MLflow untuk \textit{tracking} eksperimen. &
Reprodusibilitas meningkat melalui kontainerisasi dan pengelolaan eksperimen terpusat, sementara siklus pengembangan \textit{pipeline} menjadi lebih terstruktur. &
\autocite{amershi2019software,pimentel2019notebooks} \\
\hline
\textit{Deployment} Model &
Proses \textit{deployment} manual dengan \textit{refactoring} besar dan sering terjadi \textit{bug} pasca-\textit{deploy}. &
KServe dengan \textit{InferenceService} deklaratif dan \textit{deployment} otomatis melalui GitOps di atas Knative Serving. OpenTelemetry mengumpulkan log dan \textit{trace} inferensi ke Loki dan Grafana Tempo sebagai dasar \textit{monitoring} dan analitik lanjutan. &
Waktu \textit{deployment} berpotensi turun secara signifikan, dan insiden terkait kesalahan konfigurasi \textit{deployment} dapat dikurangi lewat otomasi. &
\autocite{paleyes2022challenges} \\
\hline
Pola \textit{Deployment} Lanjut &
Pola seperti \textit{canary release} jarang digunakan, organisasi umumnya mengandalkan \textit{binary switchover}. &
Dukungan \textit{canary release}, A/B \textit{testing}, dan \textit{automatic rollback} prediktif. &
\textit{Rollout} dapat dilakukan secara bertahap dan aman, sehingga MTTR menurun karena kegagalan dapat dibatasi pada \textit{canary traffic}. &
\autocite{shankar2022operationalizing} \\
\hline
Deteksi \textit{Drift} &
Deteksi \textit{drift} bersifat reaktif dengan jeda deteksi yang panjang terhadap isu penting dan perubahan gradual. &
\textit{Monitoring} proaktif dengan uji statistik, \textit{alert real-time}, dan penetapan ambang adaptif. &
Jeda deteksi \textit{drift} dapat dipersingkat sehingga \textit{retraining} dapat dijadwalkan lebih tepat waktu untuk menjaga kinerja model. &
\autocite{cespedes2024efficient,bayram2022concept} \\
\hline
\textit{Governance} \& \textit{Lineage} &
Dokumentasi dilakukan secara manual dan \textit{audit trail} sering tidak lengkap atau tersebar. &
DataHub dengan \textit{metadata ingestion} otomatis dan \textit{lineage} yang dapat ditelusuri dari \textit{source} hingga model, dibantu OpenLineage sebagai protokol pengirim peristiwa lintas \textit{engine} pemrosesan. &
Waktu persiapan audit dapat berkurang dan cakupan metadata meningkat melalui proses \textit{ingestion} dan \textit{lineage} otomatis. &
\autocite{stoyanovich2020responsible,lamer2024data} \\
\hline
\textit{Observability} &
Fokus \textit{monitoring} pada metrik infrastruktur, tanpa cakupan metrik ML spesifik di sebagian besar organisasi. &
\textit{Observability} mencakup metrik ML, \textit{dashboard} kinerja model, dan \textit{alert} berbasis indikator \textit{drift}. Prometheus dan Grafana menjadi inti, OpenTelemetry mengalirkan log dan \textit{trace} ke Loki dan Grafana Tempo, sementara Alertmanager mengonsolidasikan \textit{alert} ke berbagai kanal notifikasi. &
Proses \textit{debugging} dan \textit{root cause analysis} menjadi lebih efektif sehingga MTTR dapat diturunkan. &
\autocite{shankar2022operationalizing} \\
\hline
Integrasi CI/CD &
CI/CD untuk \textit{pipeline} dan model dilakukan secara manual atau semiotomatis, dengan \textit{workflow} yang tidak seragam. &
GitOps dengan Argo CD sebagai mekanisme \textit{continuous deployment} deklaratif, dilengkapi Tekton untuk \textit{continuous integration} dan Argo Rollouts untuk \textit{progressive delivery}, dengan Git sebagai sumber kebenaran tunggal yang mencakup seluruh konfigurasi DataOps dan MLOps. &
Frekuensi \textit{deployment} dapat meningkat dan variasi kesalahan \textit{deployment} dapat ditekan melalui otomasi dan \textit{peer review} pada \textit{pull request}. &
\autocite{kreuzberger2023mlops} \\
\hline
Pengalaman \textit{Developer} &
\textit{Developer} menghabiskan banyak waktu untuk tugas repetitif, \textit{onboarding} lama, dan penggunaan alat yang terpisah. &
Platform terintegrasi dengan \textit{self-service} untuk data, fitur, \textit{pipeline}, dan akses; dex sebagai \textit{identity provider} OIDC dan oauth2-proxy menyatukan otentikasi tunggal di seluruh antarmuka pengembang. &
\textit{Onboarding} dapat dipercepat dan produktivitas \textit{data scientist} meningkat karena hambatan akses infrastruktur berkurang. &
\autocite{paleyes2022challenges,rella2022mlops} \\
\hline
Efisiensi Biaya &
Utilisasi sumber daya tidak optimal dan tidak ada atribusi biaya yang jelas ke \textit{workload}. &
\textit{Auto-scaling} berbasis HPA, VPA, dan KEDA, pemantauan biaya melalui OpenCost, serta alokasi sumber daya yang lebih terukur pada klaster Kubernetes. Seluruh komponen inti berbasis \textit{open source} sehingga biaya lisensi dapat ditekan. &
Penghematan biaya signifikan dapat dicapai melalui optimasi pemakaian sumber daya dan \textit{auto-scaling}, sebagaimana ditunjukkan dalam berbagai studi kasus adopsi Kubernetes \autocite{lakka2025kubernetes}; desain ini mengadopsi prinsip yang sama untuk mengoptimalkan pemakaian sumber daya. &
\citereset\autocite{lakka2025kubernetes} \\
\hline
Keamanan \& \textit{Compliance} &
Kontrol akses beragam di tiap sistem, dan \textit{audit} tersebar tanpa pandangan terpadu. &
RBAC terintegrasi, \textit{audit trail} komprehensif, dan pengecekan \textit{compliance} yang terotomatisasi. dex sebagai \textit{identity provider} OIDC memusatkan identitas, oauth2-proxy menegakkan otentikasi pada antarmuka aplikasi, Kyverno menegakkan kebijakan admisi pada klaster, OpenBao mengelola \textit{secrets}, sementara Falco dan Trivy Operator memantau ancaman \textit{runtime} dan kerentanan citra. &
Risiko insiden keamanan dapat berkurang melalui kontrol terpusat dan proses audit menjadi lebih ringkas. &
\autocite{ahmad2024mlops} \\
\hline
\end{longtable}
}